% Generated draft: requires mathematical and source review.
\documentclass[11pt]{article}
\usepackage[margin=25mm]{geometry}
\usepackage{fontspec}
\setmainfont{Latin Modern Roman}
\usepackage{amsmath,amssymb,mathtools,unicode-math}
\setmathfont{Latin Modern Math}
\usepackage{enumitem,needspace,xurl,hyperref}
\setlength{\parindent}{0pt}
\setlength{\parskip}{5pt}
\setlength{\emergencystretch}{3em}
\newcommand{\N}{\mathbb{N}}
\newcommand{\Z}{\mathbb{Z}}
\newcommand{\Q}{\mathbb{Q}}
\newcommand{\R}{\mathbb{R}}
\newcommand{\C}{\mathbb{C}}
\newcommand{\F}{\mathbb{F}}
\newcommand{\A}{\mathbb{A}}
\newcommand{\PP}{\mathbb{P}}
\newcommand{\E}{\mathbb{E}}
\newcommand{\eps}{\varepsilon}
\newcommand{\dd}{\,\mathrm{d}}
\begin{document}
\section*{UnsolvedMath ID: 20001000}
\textbf{Source problem code:} catalogue entry
\subsection{Definitions and mathematical statement}
Leonardo Rolla's exact source question asks whether an i.i.d. chip configuration of sufficiently small density \(\mu\) on \(\Z^d\) is stabilizable with high probability. The aid reconstructs this as the ordinary fixed-energy abelian sandpile, rather than a stochastic-toppling or activated-walk model; the statement status is \texttt{reconstructed\_unverified}.

Fix \(d\ge1\), and let \(\eta(x)\), \(x\in\Z^d\), be i.i.d. nonnegative integer heights with finite mean \(\mu=\E\eta(0)\). A site with at least \(2d\) chips may topple, losing \(2d\) and sending one to each nearest neighbor. Stable heights lie in \(\{0,\ldots,2d-1\}\). For finite boxes with wired sink, let \(u_R(x)\) be the legal stabilization odometer. Infinite-volume stabilizability means that their increasing exhaustion limit \(u(x)\) is finite at every site, equivalently a complete legal stabilization topples each site only finitely often. It need not use only finitely many topplings globally. Randomness lies in initial heights.

The aid attributes to Fey--Meester--Redig the qualitative affirmative result
\[
 \mu<d\quad\Longrightarrow\quad
 \Pr(\eta\text{ is stabilizable})=1.
\]
No assertion at \(\mu=d\) is supplied. Product ergodicity makes this translation-invariant event have probability zero or one, so ``high probability'' in the intended infinite-volume question is interpreted as almost surely. Every finite wired box stabilizes regardless of density, so its stabilization probability alone does not express the question.

The aid's additional localization certificate assumes an exponential moment and some \(t>0\) satisfying
\[
 q(t)=4d^2e^{-t}\E e^{t\eta(0)}<1.
\]
It reports stabilization and the bound
\[
 \Pr(|\mathcal C_0|\ge n)\le
 \frac1{4d^2}\frac{q(t)^n}{1-q(t)}\qquad(n\ge1),
\]
where \(\mathcal C_0\) is the nearest-neighbor connected component of sites with positive odometer containing the origin, or the empty set if the origin never topples. The exact finite-volume-to-infinite-volume and cluster conventions require verification from the missing proof artifact.

For Poisson heights of mean \(0<\mu<1\), the admissible optimum is \(t=\log(1/\mu)>0\), giving
\[
 q_d(\mu)=4d^2\mu e^{1-\mu}.
\]
Thus \(q_d(\mu)<1\) in this small-mean range is the reported quantitative certificate, and \(\mu<1/(4ed^2)\) is a simpler sufficient condition. At \(\mu=0\) the configuration is identically zero. The restriction \(\mu<1\) is necessary for that optimizer to satisfy \(t>0\): for \(\mu\ge1\), the expression \(4d^2\mu e^{1-\mu}\) is not the infimum over positive \(t\), and cannot be used as a large-density certificate. This admissibility qualification is absent from the supplied summary and is explicit here. The localization criterion is sufficient and substantially narrower than the reported mean-only qualitative result; neither is asserted to characterize the sharp stabilization threshold.
\subsection{Short English statement}
Under the explicitly reconstructed ordinary sandpile model, the supplied aid reports almost-sure stabilization for i.i.d. mean below \(d\). It separately gives an exponential cluster-tail certificate under an exponential-moment condition, with the Poisson optimization restricted to its valid small-mean range.
\subsection{Sources}
\begin{itemize}
\item Local UnsolvedMath registry, permanent ID \(20001000\), source code \texttt{catalogue entry}; \texttt{aim-combinatorics-notes.json}, \texttt{notes[124]}, tag \texttt{problem}.
\item AIM workshop ``Generalizations of chip-firing and the critical group,'' Problem 35, Leonardo Rolla: \url{https://aimath.org/pastworkshops/chipfiringproblems.pdf}. Statement status \texttt{reconstructed\_unverified}. No PDF was accessed for this draft.
\item Fey--Meester--Redig, named as the source of the infinite-volume ordinary-sandpile result; no full title, date, theorem number, or URL is supplied. The exponential-tail certificate is a separate candidate result without an authoritative external citation.
\item Research attempt 1: \texttt{partial\_result}, problem status at run \texttt{solved}, rigor label \texttt{proved}, novelty confidence low; registry status \texttt{solved}, classification \texttt{SOLVED-IN-LITERATURE}. Created and updated 14 August 2026. No original remarks or literature field are supplied. This machine-generated aid is not peer reviewed; independently verify before citation.
\item Ulam AI / UnsolvedMath Contributors; CC BY 4.0.
\end{itemize}
\end{document}
